\clearpage
\section{R17: native evidence and return to the original}
\label{sec:r17-evidence}

The final R17 implementation was built with CUDA 12.8.61 and MSVC 19.44
for \code{sm\_120}. Measurements use the RTX 5070 Ti Laptop GPU, with
46 SMs, 36 MiB L2, and 12,820,480,000 reported VRAM bytes. The host is a
Core Ultra 7 255HX. GPU jobs were serialized. All R17 trials use the
same final executable; its SHA-256 is recorded in the measurement and
compiled-code reports.

\subsection{Independent correctness and compiled execution}

All nine native test suites pass, including the inherited spatial and
feedback suites. The new fused suite passes 1,513,078 checks against the
unchanged canonical C machine and an independent per-bit ASA/NA/JK model.
Its field comparisons include 722,912 State64 words, 451,820 receipt words,
and 292,264 feedback words; other checks cover errors, epochs and API
contracts. The inherited native VM suite adds 28,794 assertions, and its
feedback suite adds 1,718 checks. All 49 inherited Python tests pass.

The fused cases cover all 16 opcodes, alias-sensitive raw fields, multiple
lanes, both fetch paths and feedback bindings, four word profiles, key
hits and misses, forced two-cell banks, halted and faulted lanes, stale
EMIT flags, nonzero starting epochs, an ordered nondefault stream, and
chunks of 1, 7, 32 and 256 across 263 requested epochs. LSYS cases cover
every shift from 0 through 30, clamped arguments, negative odd values and
signed extrema, including $-2^{31}$. This is finite-machine evidence;
the exact transition proof supplies the general composition argument.

The compiled audit inspects actual SASS and resource declarations,
binding their bytes to the executable and native source. Integer texture
variants contain actual \code{TLD} sites. All six fused variants have
zero declared stack and local-memory bytes, with no \code{LDL}/\code{STL}
instructions in the inspected kernels. The final register counts and
comparison with the preserved initial build are recorded in
\path{review/R17_NATIVE_CODE.md}. These facts concern this compiler target;
they do not establish occupancy, cache hits or performance on another GPU.

The initial LSYS division was integer-valued but compiled through a
reciprocal-and-correction sequence. The final unsigned magnitude shift
and sign restoration removes that need while retaining signed truncation
toward zero. Some floating-unit instructions can still materialize constant
bit patterns in compiler output. Exact integer state semantics must not be
confused with a claim that every generated opcode uses an integer unit.

\subsection{Matched device and complete-host measurements}

The tiny workload is the original three-cell source fixture, with 96
requested epochs and identical initial lane words. The larger workload
contains 1,048,576 synthetic Cell48 instructions: 48 MiB of operator cells,
plus the 128-byte header. Its 15 non-HALT opcodes, arguments and successor
links are generated from an explicit recipe; lane initialization spreads
raw rho, lineage and starting cell. It runs 128 epochs on 65,536 lanes.
The larger program is generated input, rather than precomputed VM output.

Each case compares R17 reference/fused dispatch and texture/global reads.
The tiny cases additionally run the unchanged R16 executable. Three fresh
process trials per variant give 132 total runs, with rotated ordering and
no warmup. Every case has matching full-array FNV-1a-64 and exact first-lane
fields, with the expected accepted EMIT totals in the tiny coupled cases.
FNV is noncryptographic; the independent tests separately compare complete
individual fields. Chunk size is 32 throughout the measured fused runs.

\begin{center}\small
\begin{tabular}{@{}l r l r r r r@{}}
\toprule
Program & Lanes & Mode & Ref. tex. & Fused tex. & Tex. gain & Global gain\\
 & & & \multicolumn{2}{c}{Device median (ms)} & \multicolumn{2}{c}{Reference/fused}\\\midrule
3 cells & 1 & pure & 0.803 & 0.082 & 9.76$\times$ & 7.78$\times$\\
3 cells & 1 & copy32 & 2.337 & 0.295 & 7.92$\times$ & 7.01$\times$\\
3 cells & 4,096 & pure & 0.749 & 0.064 & 11.66$\times$ & 14.92$\times$\\
3 cells & 4,096 & copy32 & 2.387 & 0.212 & 11.25$\times$ & 11.35$\times$\\
3 cells & 262,144 & pure & 5.992 & 0.336 & 17.82$\times$ & 20.28$\times$\\
3 cells & 262,144 & copy32 & 16.051 & 0.830 & 19.33$\times$ & 20.05$\times$\\
48 MiB & 65,536 & pure & 3.003 & 0.897 & 3.35$\times$ & 4.21$\times$\\
48 MiB & 65,536 & copy32 & 5.767 & 1.135 & 5.08$\times$ & 6.21$\times$\\
\bottomrule
\end{tabular}\end{center}


The table reports medians of device batch time in milliseconds. Across
these 16 matched R17 comparisons, device ratios range from
$3.347\times$ to $20.277\times$. The unchanged R16 tiny-fixture comparisons
give $6.460\times$ to $19.710\times$. Setup, final readback and host work
are outside device batch time. These are measurements of the stated word
machine workloads, with three repeats, rather than a universal speed bound.

Cold host-complete ratios range from $0.965\times$ to $1.098\times$:
13 of 16 cases have a lower fused median, and three have a higher median.
Initialization, allocation, upload, final readback and hashing remain
substantial costs. For example, a large synthetic table has a different
setup cost from the three-cell fixture. Device speed ratios must not be
presented as complete application speedups.

The texture path is not uniformly faster than global reads; the raw trials
retain both results. A 48 MiB program exceeds the reported L2 capacity,
but that alone does not establish that every cell is visited, that a cache
is saturated, or that the device reaches a memory-bandwidth limit. No new
S2 operation is measured here. Earlier S2 results retain their original
inputs and correctness contracts.

\subsection{What returning to the original source changes}

The fresh audit returns to the user's TOM World Query Kernel 0.6.0
directory, preserving its files. Two concrete mechanisms deserve further
integration. First, the original GPU sources express several phase and
seam operations with exact 32-bit modular arithmetic. Their input ranges,
floor-division behavior and overflow treatment must be retained when
replacing wider intermediate arithmetic. Second, the original seeded
artifact workflow consumes a complete ordered stream of executed EMIT
records. A final receipt is insufficient to reconstruct that stream.

The seeded compiler validates dependencies and evaluates its formal tree
on the host, then encodes the result and lowers literal bytes through
\code{emit.graph}. R17's native transition port does not yet make that
formal evaluator a GPU program. A useful next integration is an opt-in,
lossless EMIT sink with explicit capacity and ordering, verified against
the original seeded compiler and materializer. It must retain every fresh
EMIT inside a fused chunk, including repeated payloads and EMIT-HALT.
Its additional output work requires separate matched measurements.

The original packed key and Morton notation are also distinct. RADIX
reads the canonical packed key, while lineage incorporates canonical
cell identity. A derived cache may rearrange physical storage only while
preserving those logical keys, indices and successor semantics. Reordering
the canonical program by opcode or Morton code would change execution.

\path{review/R17_ORIGINAL_RETURN.md} and
\path{review/R17_SEEDED_RUNTIME_AUDIT.md} identify the source locations and
bounded next steps. The reports distinguish reused mechanisms, verified
R17 behavior and capabilities still absent. The source recipe, trials,
compiled audit, checks and mathematical proof remain independently
inspectable in the release.

\clearpage
